\markdownRendererDocumentBegin
To achieve this, we need to modify \texttt{1cyl.par} once more, updating the initial condition (\textit{startFrom}) to the computed base flow, setting the integration time (\textit{endTime}), and configuring \textit{userParam} for the direct eigenproblem mode:\markdownRendererInterblockSeparator
{}\markdownRendererInputFencedCode{./_markdown_main/ac766bb5ecff9d66923aff56380e4230.verbatim}{bash}\markdownRendererInterblockSeparator
{}\hline \bigskip As in the previous section, you can execute the code using the following command:\markdownRendererInterblockSeparator
{}\markdownRendererInputFencedCode{./_markdown_main/be436a55e82dbfde012ae4427ee71a63.verbatim}{bash}\markdownRendererInterblockSeparator
{}To monitor the progress of the output, use:\markdownRendererInterblockSeparator
{}\markdownRendererInputFencedCode{./_markdown_main/1c16059fb9c9da4cacf53874807afcd6.verbatim}{bash}\markdownRendererInterblockSeparator
{}To exit, press ``Ctrl+c''.\markdownRendererInterblockSeparator
{}Upon completion, you can plot both spectrum (shown in figure \ref{fig:spectra}):\markdownRendererInterblockSeparator
{}\markdownRendererInputFencedCode{./_markdown_main/a167b6b578db6353b2ca9d5b1020b61b.verbatim}{bash}\markdownRendererDocumentEnd